\begingroup\small
\begin{longtable}{l p{7.2cm} l}
\hline
statement & title & type \\
\hline
\endhead
Lemma~III.1.1 & Bajpai, Bennett and Chan {\cite[proof of Theorem~1.1]{BajpaiBennettChan2024}} & literature \\
Proposition~III.1.2 & the head of the gap spectrum & computed \\
Remark~III.1.3 & a ceiling for common divisors & computed \\
Remark~III.1.4 & a counting model, measured & measured \\
Remark~III.1.5 & common factors of pairs & measured \\
Remark~III.2.1 & the values $x^2 + 1$ & computed \\
Proposition~III.2.2 & two explicit chains of powerful values of $\Kreis{3}$ and $\Kreis{6}$ & proved \\
Proposition~III.2.3 & powerful values up to a bound & computed \\
Remark~III.2.4 & a class without squares & computed \\
Proposition~III.2.5 & the simplest case of $\Kreis{8}$ & proved \\
Remark~III.3.1 & no joint event along the divisors, and the power of the test & measured \\
Remark~III.4.1 & no deviation from equidistribution, memory or coupling & measured \\
Question~III.5.1 & the most frequent gap & open \\
Question~III.5.2 & powerful values of $x^4 + 1$ & open \\
Remark~III.5.3 & a conjecture on the most frequent gap & heuristic \\
\hline
\end{longtable}
\endgroup

The types: \emph{proved}, proved in this part; the proof may use published results (2); \emph{literature}, due to the cited source, sometimes with a proof given here for completeness (1); \emph{computed}, the result of an exact computation, or a proof that uses one of our computations (5); \emph{measured}, a measurement, with a model that is not proved (4); \emph{heuristic}, a heuristic, used as a yardstick only (1); \emph{open}, an open question (2).
